\begin{proposition}[Tangent index and Picard components]\hypertarget{tag-01BN}{}\label{cdp:st:C.1.1}
\StatementRecord{cdp:st:C.1.1}{proved}{01BN}{the definitions and constructions stated in the proof}

Let $Y$ be a connected smooth compact complex threefold with
$b_1(Y)=b_2(Y)=0$. Every holomorphic line bundle $L$ on $Y$ satisfies
\[
 \chi(Y,T_Y\otimes L)=\frac12\int_Yc_3(T_Y)
 =1-\frac{b_3(Y)}2.
\]
The orientation in the integral is the complex orientation. The group
$\Pic(Y)$ has finitely many connected components, each a torsor for the
complex vector group $H^1(Y,\OO_Y)$, and
$\pi_0\Pic(Y)\simeq H^2(Y;\Z)$ through the first Chern class.
\end{proposition}


\begin{proof}

\begin{enumerate}[label={$\langle1\rangle$\arabic*.},leftmargin=*]
\item The universal coefficient theorem gives
$H^1(Y;\Z)=\Hom(H_1(Y;\Z),\Z)=0$ from $b_1(Y)=0$.
The group $H^2(Y;\Z)$ is finite because $Y$ has a finite CW model and
$b_2(Y)=0$. In the exponential sequence, normalized by
$\exp(2\pi i\,\cdot)$, the relevant maps are
\[
0=H^1(Y;\Z)\longrightarrow H^1(Y,\OO_Y)
 \xrightarrow{\exp} \Pic(Y)
 \xrightarrow{c_1} H^2(Y;\Z).
\]
They are group homomorphisms; the middle map is holomorphic on the
Picard group. Consequently its image $\Pic^0(Y)$ is isomorphic to the
complex vector space $H^1(Y,\OO_Y)$. The next map in the exponential
sequence is $H^2(Y;\Z)\to H^2(Y,\OO_Y)$; it is zero because its
source is finite and its target is a complex vector space. Thus
$c_1$ is surjective, and the component group is exactly $H^2(Y;\Z)$.

\item Poincar\'e duality gives $H^4(Y;\Q)=0$. Thus
$c_1(T_Y)$, $c_2(T_Y)$ and $c_1(L)$ vanish in rational cohomology.
The degree-six component of
$\operatorname{ch}(T_Y)\exp(c_1(L))\operatorname{td}(T_Y)$ is then
$c_3(T_Y)/2$. The Hirzebruch--Riemann--Roch formula for the Dolbeault
index gives the first equality. Chern--Gauss--Bonnet identifies
$\int_Yc_3(T_Y)$ with $e(Y)$. Since $Y$ is a connected closed oriented
six-manifold, its Betti numbers satisfy
$b_0=b_6=1$, $b_1=b_5=0$, and $b_2=b_4=0$. Hence
$e(Y)=2-b_3(Y)$. The pairing on $H^3(Y;\Q)$ is nondegenerate and
alternating, so $b_3(Y)$ is even, consistently with the integral index.


\end{enumerate}
\end{proof}